% Quelle: 6. Übung CT Lösungen.pdf
% Art: Übung mit Lösungen
% Themen: Kreisteilungsklassen, Generatorpolynom zyklischer Codes, Kontroll- und Generatormatrix in Standardform, Syndrom, Inverse modulo 11, Syndrom-Decodierung [10,8,3]_11-Code
% Hinweis: Die Quelle ist ein Übungsblatt mit Lösungen; hier stehen nur die in den Lösungen verwendeten allgemeinen Formeln/Hilfsmittel (wörtlich aus der Quelle). Aufgaben und Lösungswege: latex/uebungen/06_uebung.tex.

\section{Übung 6: Wissen aus den Lösungen}
\subsection{Zyklische Codes}

\begin{bemerkung}[Kreisteilungsklassen, aus 6\_1 b) und 6\_5]
Kreisteilungsklassen: $\{s, sp, sp^2, \dots, sp^{d-1}\}$
\end{bemerkung}

\begin{bemerkung}[Generatorpolynom, aus 6\_2 und 6\_3]
Generatorpolynom $g(X) \mid X^n - 1$, \quad $\operatorname{grad} g(X) = p = n - k$.
% In der Quelle konkret für n = 7: g(X) | X^7 - 1.
\end{bemerkung}

\subsection{Kontroll- und Generatormatrix, Syndrom}

\begin{bemerkung}[aus 6\_4 und 6\_8]
$H = (\, E \mid A \,)$, \qquad $G = (-A^T \mid E\,)$\\
$\operatorname{Syn}(\mathbf{r}) = H\mathbf{r}^T$
\end{bemerkung}

\subsection{$[10, 8, 3]_{11}$ Code}

\begin{bemerkung}[Inverse Modulo 11, aus 6\_9]
\begin{center}
\begin{tabular}{c|rrrrrrrrrr}
$a$      & 1 & 2  & 3 & 4 & 5  & $-5$ & $-4$ & $-3$ & $-2$ & $-1$\\
\midrule
$a^{-1}$ & 1 & $-5$ & 4 & 3 & $-2$ & 2 & $-3$ & $-4$ & 5 & $-1$\\
\end{tabular}
\end{center}
\end{bemerkung}

\begin{algorithmus}[{Decodierung des $[10, 8, 3]_{11}$ Codes, aus 6\_9}]
\[
H = \begin{pmatrix}
1 & 1 & 1 & 1 & 1 & 1 & 1 & 1 & 1 & 1\\
1 & 2 & 3 & 4 & 5 & 6 & 7 & 8 & 9 & 10
\end{pmatrix}
\equiv
\begin{pmatrix}
1 & 1 & 1 & 1 & 1 & 1 & 1 & 1 & 1 & 1\\
1 & 2 & 3 & 4 & 5 & -5 & -4 & -3 & -2 & -1
\end{pmatrix}
\]
$\mathbf{e} = \mathbf{r} - \mathbf{c} = (\,0 \dots m \dots 0\,)$, \quad
$S(\mathbf{r}) = \begin{pmatrix} s_0\\ s_1 \end{pmatrix} = \begin{pmatrix} 1\\ x \end{pmatrix} y$, \quad
$\mathbf{c} = \mathbf{r} - \mathbf{e}$
\end{algorithmus}
